
\documentclass{article}
\usepackage{colortbl}
\usepackage{makecell}
\usepackage{multirow}
\usepackage{supertabular}

\begin{document}

\newcounter{utterance}

\centering \large Interaction Transcript for game `codenames', experiment `frequency\_low', episode 3 with reflexion{-}llama{-}3.1{-}8b{-}instruct.
\vspace{24pt}

{ \footnotesize  \setcounter{utterance}{1}
\setlength{\tabcolsep}{0pt}
\begin{supertabular}{c@{$\;$}|p{.15\linewidth}@{}p{.15\linewidth}p{.15\linewidth}p{.15\linewidth}p{.15\linewidth}p{.15\linewidth}}
    \# & $\;$A & \multicolumn{4}{c}{Game Master} & $\;\:$B\\
    \hline

    \theutterance \stepcounter{utterance}  
    & & & \multicolumn{2}{p{0.3\linewidth}}{
        \cellcolor[rgb]{0.9,0.9,0.9}{
            \makecell[{{p{\linewidth}}}]{
                \texttt{\tiny{[GM$|$GM]}}
                {'hidden': {'team': ['eggroll', 'plantain', 'unchallenging'], 'innocent': ['bandmaster', 'ladylove', 'customizer', 'defenseman', 'narrowness', 'neckband'], 'opponent': [], 'assassin': ['seacoast']}, 'revealed': {'team': {'team': [], 'innocent': [], 'opponent': [], 'assassin': []}, 'opponent': {'team': [], 'innocent': [], 'opponent': [], 'assassin': []}}}
            }
        }
    }
    & & \\ \\

    \theutterance \stepcounter{utterance}  
    & & \multicolumn{4}{p{0.6\linewidth}}{
        \cellcolor[rgb]{0.9,0.9,0.9}{
            \makecell[{{p{\linewidth}}}]{
                \texttt{\tiny{[P1$\langle$GM]}}
                \texttt{You have attempted this same task before and did not succeed. These are the plans you wrote after those attempts.} \\
\texttt{Trial \#1: New plan: } \\
\\ 
\texttt{In my previous attempts, I was able to guess some of the team words correctly, but I failed to guess all of them before the opposing team uncovered all of theirs. Upon reflection, I realize that I was focusing too much on individual words and not enough on the overall strategy of the game. } \\
\\ 
\texttt{In the game, the goal is to have your partner guess all of your team's words before the opposing team has uncovered all of theirs. To achieve this, I need to think about the relationships between the words and how they can be grouped together to make it easier for my partner to guess them.} \\
\\ 
\texttt{In my previous attempts, I was guessing words that were related to each other, but I was not considering the broader categories that these words belonged to. For example, when I guessed the word "typewrite", I was thinking about the word "writing", but I did not consider the category of "office equipment" that both words belong to.} \\
\\ 
\texttt{To improve my strategy, I will focus on identifying the broader categories that the words belong to and use these categories to guide my guesses. I will also try to group the words together in a way that makes it easier for my partner to guess them. For example, if I have a word that is related to "food", I will try to guess other words that are also related to "food" and use these words to help my partner guess the remaining words.} \\
\\ 
\texttt{Additionally, I will pay closer attention to the words that are being guessed by the opposing team and try to use this information to my advantage. If I see that the opposing team is guessing words that are related to a particular category, I will try to guess words that are also related to that category, but are not as obvious.} \\
\\ 
\texttt{By focusing on the broader categories and relationships between the words, I hope to improve my chances of winning the game.} \\
\\ 
\texttt{Trial \#2: New plan: In my previous attempts, I was focusing too much on individual words and not enough on the overall strategy of the game. I realized that I need to think about the relationships between the words and how they can be grouped together to make it easier for my partner to guess them. } \\
\\ 
\texttt{To improve my strategy, I will focus on identifying the broader categories that the words belong to and use these categories to guide my guesses. I will also try to group the words together in a way that makes it easier for my partner to guess them. } \\
\\ 
\texttt{In this specific game, I noticed that my team's words "eggroll" and "plantain" are both related to food. However, I failed to guess the third word "unchallenging" which is not related to food. } \\
\\ 
\texttt{Next time, I will try to identify the broader categories that all three words belong to. Since "eggroll" and "plantain" are related to food, I will try to find a category that includes all three words. I will also pay closer attention to the words that are being guessed by the opposing team and try to use this information to my advantage. } \\
\\ 
\texttt{For example, if I see that the opposing team is guessing words that are related to a particular category, I will try to guess words that are also related to that category, but are not as obvious.} \\
\\ 
\\ 
\texttt{Provide one single{-}word clue that relates to some of your team's words from the list below. You can choose to how many and to which words your clue relates to.} \\
\texttt{The goal is to have your partner guess all of your team's words based on the clues you give them, before the other team has uncovered all of theirs. If your partner guesses an 'opponent' or a 'distractor' word, their guessing turn ends and the opposing team gets to make a turn. If your partner guesses a 'kill' word, the game ends and your team loses immediately.} \\
\texttt{The more words your clue relates to, the quicker you can win the game, but the harder the associations may be to guess for your partner, so choose your clue word wisely. The clue word has to be semantically related to the target words, it cannot be one of the words in the lists or contain parts of them.} \\
\\ 
\texttt{Always give your single{-}word clue and your comma{-}separated list of related target words in the following format and make your answers as short as possible, never include any other text than is required in this form:} \\
\\ 
\texttt{CLUE: <WORD>} \\
\texttt{TARGETS: <TARGETS>} \\
\\ 
\texttt{Your team words are: eggroll, plantain, unchallenging.} \\
\texttt{Your opponent's team words are: .} \\
\texttt{Distractor words are: bandmaster, ladylove, customizer, defenseman, narrowness, neckband.} \\
\texttt{The kill words are: seacoast.} \\
            }
        }
    }
    & \\ \\

    \theutterance \stepcounter{utterance}  
    & \multicolumn{4}{p{0.6\linewidth}}{
        \cellcolor[rgb]{0.9,0.9,0.9}{
            \makecell[{{p{\linewidth}}}]{
                \texttt{\tiny{[P1$\rangle$GM]}}
                \texttt{CLUE: cuisine} \\
\texttt{TARGETS: eggroll, plantain} \\
            }
        }
    }
    & & \\ \\

    \theutterance \stepcounter{utterance}  
    & & & \multicolumn{2}{p{0.3\linewidth}}{
        \cellcolor[rgb]{0.9,0.9,0.9}{
            \makecell[{{p{\linewidth}}}]{
                \texttt{\tiny{[GM$|$GM]}}
                \texttt{cuisine} \\
            }
        }
    }
    & & \\ \\

    \theutterance \stepcounter{utterance}  
    & & & \multicolumn{2}{p{0.3\linewidth}}{
        \cellcolor[rgb]{0.9,0.9,0.9}{
            \makecell[{{p{\linewidth}}}]{
                \texttt{\tiny{[GM$|$GM]}}
                ['eggroll', 'plantain']
            }
        }
    }
    & & \\ \\

    \theutterance \stepcounter{utterance}  
    & & & \multicolumn{2}{p{0.3\linewidth}}{
        \cellcolor[rgb]{0.9,0.9,0.9}{
            \makecell[{{p{\linewidth}}}]{
                \texttt{\tiny{[GM$|$GM]}}
                {'word': 'eggroll', 'assignment': 'team'}
            }
        }
    }
    & & \\ \\

    \theutterance \stepcounter{utterance}  
    & & & \multicolumn{2}{p{0.3\linewidth}}{
        \cellcolor[rgb]{0.9,0.9,0.9}{
            \makecell[{{p{\linewidth}}}]{
                \texttt{\tiny{[GM$|$GM]}}
                {'word': 'plantain', 'assignment': 'team'}
            }
        }
    }
    & & \\ \\

    \theutterance \stepcounter{utterance}  
    & & & \multicolumn{2}{p{0.3\linewidth}}{
        \cellcolor[rgb]{0.9,0.9,0.9}{
            \makecell[{{p{\linewidth}}}]{
                \texttt{\tiny{[GM$|$GM]}}
                \texttt{opponent has won} \\
            }
        }
    }
    & & \\ \\

    \theutterance \stepcounter{utterance}  
    & \multicolumn{4}{p{0.6\linewidth}}{
        \cellcolor[rgb]{0.95,0.95,0.95}{
            \makecell[{{p{\linewidth}}}]{
                \texttt{\tiny{[P2$\rightarrow$P1]}}
                \texttt{Post{-}episode Reflexion (not sent to GM):} \\
\texttt{New plan: In my previous attempts, I was focusing too much on individual words and not enough on the overall strategy of the game. I realized that I need to think about the relationships between the words and how they can be grouped together to make it easier for my partner to guess them. } \\
\\ 
\texttt{To improve my strategy, I will focus on identifying the broader categories that the words belong to and use these categories to guide my guesses. I will also try to group the words together in a way that makes it easier for my partner to guess them. } \\
\\ 
\texttt{In this specific game, I noticed that my team's words "eggroll" and "plantain" are both related to food, and "unchallenging" is related to a concept or a state. However, I failed to find a category that includes all three words. } \\
\\ 
\texttt{Next time, I will try to identify a category that is common to all three words. For example, I will think about the context in which these words are used. Are they related to a particular activity, a type of experience, or a state of being? By finding a category that includes all three words, I will be able to make more informed guesses and increase my chances of winning the game.} \\
            }
        }
    }
    & & \\ \\

\end{supertabular}
}

\end{document}
